%-----------------------------------
% Erklaerung zur Nutzung von KI-Hilfsmitteln
% Steht hinter dem KI-Verzeichnis, ohne Gliederungsnummer, im Inhaltsverzeichnis.
% Fassung laut Vorlage "KI-Verzeichnis_fertig.pdf" des Verfassers (22.11.2026).
%-----------------------------------
\newpage
\phantomsection
\addcontentsline{toc}{section}{Erkl\"arung zur Nutzung von KI-Hilfsmitteln}
\section*{Erkl\"arung zur Nutzung von KI-Hilfsmitteln}

Im Rahmen der Erstellung dieser Arbeit wurden KI-basierte Systeme zur Unterst\"utzung des Forschungs- und Schreibprozesses genutzt. Dabei wurden die Vorgaben der FOM Hochschule f\"ur Oekonomie \& Management gem\"a{\ss} dem Leitfaden zur Gestaltung wissenschaftlicher Arbeiten (ITM/ING, M\"arz 2024, Version~1.4), die Leitlinien der Deutschen Forschungsgemeinschaft zur Sicherung guter wissenschaftlicher Praxis (Kodex~2023) sowie die Datenschutz-Grundverordnung beachtet.

\subsection*{Anwendungsbereiche}

\begin{itemize}[leftmargin=1.5em]
	\item \textbf{Themenfindung und Gliederung:} Entwicklung und Bewertung von Titelvarianten, Forschungsfragen und Inhaltsverzeichnis mit Claude Opus~5.5
	\item \textbf{Literaturarbeit:} Entwicklung von Suchbegriffen und Suchstrings sowie Quellensuche mit Claude Fable~5.0 und Claude Opus~5.5; Pr\"ufung der bibliografischen Angaben im Literaturverzeichnis mit Claude Opus~5.5
	\item \textbf{Feedback zur Arbeit:} kapitelweise Begutachtung mehrerer Fassungen nach Argumentation, Stil, Themenfokus, Methodik und Redundanz mit Claude Opus~5.5
	\item \textbf{Textarbeit:} Formulierungsvorschl\"age, Textbausteine f\"ur einzelne Passagen (insbesondere Abschnitt~4.3.1) sowie \"Uberarbeitung einzelner S\"atze zur Beseitigung von Widerspr\"uchen und zur Vereinheitlichung von Fachbegriffen mit Claude Opus~5.5
	\item \textbf{Abbildungen:} Erstellung von Abbildung~7 auf Grundlage des eigenen Modellentwurfs sowie Korrektur von Abbildung~5 mit Claude Opus~5.5
	\item \textbf{Qualit\"atspr\"ufung:} inhaltliche und formale Pr\"ufung mit Claude Opus~5.5, sprachliche Pr\"ufung mit Overleaf AI und abschlie{\ss}ende Gesamtpr\"ufung mit Claude Fable~5.2
\end{itemize}

\subsection*{Grunds\"atze der Nutzung}

\begin{itemize}[leftmargin=1.5em]
	\item Keine Eingabe oder Verarbeitung personenbezogener oder vertraulicher Daten
	\item Kritische Pr\"ufung und gegebenenfalls Anpassung aller \"ubernommenen Inhalte
	\item Kennzeichnung KI-gest\"utzt erstellter Inhalte im Text und in den Abbildungsunterschriften gem\"a{\ss} Abschnitt~6.2.4 des Leitfadens
	\item Dokumentation der Nutzung im KI-Verzeichnis und in Anhang~II
\end{itemize}

\subsection*{Eigenleistung und Dokumentation}

Alle wesentlichen inhaltlichen Entscheidungen sowie die kritische Pr\"ufung und Interpretation wurden von mir als Verfasser vorgenommen. Vorschl\"age und Inhalte der KI-Systeme habe ich nur \"ubernommen, nachdem ich sie gepr\"uft und bei Bedarf \"uberarbeitet habe. S\"amtliche Quellen wurden nach wissenschaftlichen Standards gepr\"uft und zitiert.

Die verwendeten Prompts sind in Kurzform in Anhang~II aufgef\"uhrt. Die vollst\"andigen Chatverl\"aufe sowie die Versionsst\"ande der Arbeit im Git-Repository liegen vor und k\"onnen auf Anfrage eingesehen werden.

Ich versichere, dass die vorliegende Arbeit trotz der KI-Unterst\"utzung meine eigenst\"andige wissenschaftliche Leistung darstellt.

\vspace{1.5em}
\noindent D\"usseldorf, den \rule{3.2cm}{0.4pt}

\vspace{2.5em}
\noindent\rule{6.5cm}{0.4pt}\\
Tim Wehrens
